\section{Training Objectives}

The primary training objective was to reconstruct the dynamic BOLD activity
as faithfully as possible from the compressed latent representation.
Pointwise agreement between the observed and reconstructed signals provides
a direct measure of reconstruction accuracy, but does not necessarily ensure
that the temporal evolution of the signal or the functional relationships
between brain regions are preserved. The training framework therefore
combined a primary pointwise reconstruction loss with optional auxiliary
reconstruction terms targeting temporal dynamics, static functional
connectivity, and dynamic connectivity variability. Additional variational
and supervised objectives were incorporated where applicable.

\subsection{Pointwise Reconstruction Loss}

The pointwise reconstruction loss measures the difference between the
observed BOLD signal $\mathbf{X}$ and its reconstruction
$\hat{\mathbf{X}}$. Three alternative reconstruction objectives were
considered: mean squared error (MSE), mean absolute error (MAE), and Huber
loss.

The MSE is defined as
\begin{equation}
    \mathcal{L}_{\mathrm{MSE}}
    =
    \frac{1}{N}
    \sum_{i=1}^{N}
    \left(
        X_i-\hat{X}_i
    \right)^2,
\end{equation}
where $N$ denotes the total number of reconstructed BOLD values. Squaring
the residuals places greater emphasis on larger reconstruction errors.

The MAE is defined as
\begin{equation}
    \mathcal{L}_{\mathrm{MAE}}
    =
    \frac{1}{N}
    \sum_{i=1}^{N}
    \left|
        X_i-\hat{X}_i
    \right|.
\end{equation}
Compared with MSE, MAE is less sensitive to isolated large residuals and
penalises reconstruction errors linearly.

Huber loss~\cite{Huber1964} combines quadratic and linear penalisation,
\begin{equation}
    \mathcal{L}_{\mathrm{Huber}}
    =
    \frac{1}{N}
    \sum_{i=1}^{N}
    \begin{cases}
        \frac{1}{2}(X_i-\hat{X}_i)^2,
        & |X_i-\hat{X}_i| \leq \delta, \\[4pt]
        \delta
        \left(
            |X_i-\hat{X}_i|-\frac{1}{2}\delta
        \right),
        & |X_i-\hat{X}_i| > \delta,
    \end{cases}
\end{equation}
where $\delta$ determines the transition between the quadratic and linear
regions. Huber loss therefore retains the sensitivity of MSE to small
residuals while reducing the influence of larger reconstruction errors.

These objectives were compared to determine which form of pointwise
penalisation was most suitable for reconstruction of the regional BOLD
signals. RMSE was not included as a separate training objective, as it is
a monotonic transformation of MSE and therefore provides a closely related
optimisation criterion.

\subsection{Functional Connectivity Loss}

Pointwise reconstruction alone does not explicitly constrain the
inter-regional relationships present in the BOLD signal. A functional
connectivity (FC) loss was therefore introduced to encourage preservation
of the static correlation structure between brain regions.

For an input
\[
\mathbf{X}\in\mathbb{R}^{T\times R},
\]
each regional time series is first centred across the temporal dimension
and normalised to unit Euclidean norm. The resulting FC matrix can then be
computed as
\begin{equation}
    \mathbf{C}
    =
    \bar{\mathbf{X}}^{\mathsf{T}}
    \bar{\mathbf{X}},
\end{equation}
where $\bar{\mathbf{X}}$ denotes the centred and normalised time series.
This operation is equivalent to computing the Pearson correlation between
each pair of regional signals. The reconstructed FC matrix
$\hat{\mathbf{C}}$ is calculated in the same manner from
$\hat{\mathbf{X}}$.

The FC loss is then defined as
\begin{equation}
    \mathcal{L}_{\mathrm{FC}}
    =
    \operatorname{MSE}
    \left(
        \mathbf{C},
        \hat{\mathbf{C}}
    \right).
\end{equation}

This term complements the pointwise reconstruction objective by encouraging
the reconstructed signals to preserve static inter-regional functional
organisation.

\subsection{Temporal and Dynamic Connectivity Losses}

Two additional objectives were considered to more directly constrain the
dynamic properties of the reconstructed BOLD signals: a first-derivative
loss and a sliding-window FC variability loss.

\subsubsection{First-Derivative Loss}

A reconstruction may achieve low pointwise error while remaining overly
smooth and failing to reproduce local temporal changes. To explicitly
constrain the temporal evolution of each regional signal, the first
temporal difference is computed as
\begin{equation}
    \Delta X_{t,r}
    =
    X_{t+1,r}-X_{t,r},
\end{equation}
with the equivalent quantity $\Delta\hat{X}_{t,r}$ calculated for the
reconstruction.

The derivative loss is defined as
\begin{equation}
    \mathcal{L}_{\Delta}
    =
    \operatorname{MSE}
    \left(
        \Delta\mathbf{X},
        \Delta\hat{\mathbf{X}}
    \right).
\end{equation}

This objective encourages the reconstructed BOLD signals to preserve the
magnitude and direction of changes between consecutive time points rather
than matching only their absolute values.

\subsubsection{Sliding-Window FC Variability Loss}

Static FC summarises functional relationships over the complete scan and
therefore does not capture changes in connectivity over time. To introduce
a lightweight constraint on dynamic functional connectivity, the BOLD
signals are divided into overlapping temporal windows and an FC matrix is
computed independently within each window.

For each regional pair, the standard deviation of its FC value across
windows is then calculated,
\begin{equation}
    \mathbf{S}
    =
    \operatorname{Std}_{w}
    \left(
        \mathbf{C}^{(w)}
    \right),
\end{equation}
where $\mathbf{C}^{(w)}$ denotes the FC matrix of window $w$. The
corresponding variability matrix $\hat{\mathbf{S}}$ is obtained from the
reconstructed signal.

The sliding-window FC variability loss is defined as
\begin{equation}
    \mathcal{L}_{\mathrm{SwFC-var}}
    =
    \operatorname{MSE}
    \left(
        \mathbf{S},
        \hat{\mathbf{S}}
    \right).
\end{equation}

This objective does not attempt to reproduce every individual
sliding-window FC matrix. Instead, it constrains the reconstructed signals
to reproduce the degree to which each functional connection varies over
time, providing a computationally lighter approximation to dynamic
connectivity preservation.

\subsection{Variational and Auxiliary Objectives}

The preceding losses all constrain properties of the reconstructed BOLD
signal. Additional objectives were introduced for specific model variants.

For the variational autoencoders, the reconstruction objective was
augmented with the KL-divergence term described in the Variational
Autoencoder Extension section,
\begin{equation}
    \mathcal{L}_{\mathrm{VAE}}
    =
    \mathcal{L}_{\mathrm{reconstruction}}
    +
    \beta\mathcal{L}_{\mathrm{KL}},
\end{equation}
where $\beta$ controls the strength of the latent-space regularisation.

For models containing auxiliary biomarker-prediction heads, prediction
errors were calculated using Smooth L1 loss. When multiple prediction
heads were present, their losses were averaged before being added to the
training objective,
\begin{equation}
    \mathcal{L}_{\mathrm{pred}}
    =
    \frac{1}{M}
    \sum_{m=1}^{M}
    \mathcal{L}_{\mathrm{SmoothL1}}^{(m)},
\end{equation}
where $M$ denotes the number of prediction targets.

For models containing an auxiliary diagnostic classification head,
cross-entropy loss was used,
\begin{equation}
    \mathcal{L}_{\mathrm{cls}}
    =
    -\frac{1}{B}
    \sum_{b=1}^{B}
    \log p_{b,y_b},
\end{equation}
where $p_{b,y_b}$ denotes the predicted probability assigned to the
ground-truth class $y_b$ for sample $b$.

The architectures and purposes of these auxiliary heads are described
separately in the Auxiliary Supervision section.

\subsection{Combined Objective}

The complete training framework allows the different objectives to be
combined through independently weighted loss terms. The general objective
can be expressed as
\begin{equation}
\begin{split}
    \mathcal{L}_{\mathrm{total}}
    ={}&
    \mathcal{L}_{\mathrm{recon}}
    + \lambda_{\mathrm{FC}}\mathcal{L}_{\mathrm{FC}}
    + \lambda_{\Delta}\mathcal{L}_{\Delta} \\
    &+
    \lambda_{\mathrm{SwFC}}
    \mathcal{L}_{\mathrm{SwFC-var}}
    + \beta\mathcal{L}_{\mathrm{KL}} \\
    &+
    \lambda_{\mathrm{pred}}\mathcal{L}_{\mathrm{pred}}
    + \lambda_{\mathrm{cls}}\mathcal{L}_{\mathrm{cls}} .
\end{split}
\end{equation}

This expression represents the general loss framework rather than a single
objective applied in every experiment. Terms that were not relevant to a
given model or experimental configuration were assigned a weight of zero
and therefore did not contribute to optimisation.

The pointwise, FC, derivative, and sliding-window FC variability terms
target complementary properties of reconstruction quality. Together, they
allow the optimisation objective to consider not only agreement between
individual BOLD values, but also local temporal evolution, static
inter-regional connectivity, and variability in functional connectivity
over time.